% ============================================================
% Come si usano, argomento 13: media di funzioni e varianza nel continuo
% ============================================================
\argomento{Media di funzioni e varianza nel continuo}{Formule a p.~\pageref*{f:valoremedioc}}[Media e varianza
  nel continuo]

\begin{scelta}{Dal discreto al continuo: il dizionario}
\renewcommand{\arraystretch}{1.45}
\begin{tabular}{@{}P{0.24\linewidth}P{0.33\linewidth}P{0.37\linewidth}@{}}
  \toprule
  \textbf{Concetto} & \textbf{Variabile discreta} & \textbf{Variabile continua} \\
  \midrule
  descrizione completa & pmf $p_X(x) = P(X = x)$ & densità $f_X(x)$, con $P(X = x) = 0$ \\
  probabilità di un insieme & $\sum_{x \in A} p_X(x)$ & $\int_A f_X(x) \dd x$ \\
  normalizzazione & $\sum_x p_X(x) = 1$ & $\int f_X(x) \dd x = 1$ \\
  CDF $F_X(x) = P(X \leq x)$ & a scala, salti alti $p_X(x)$ & continua, $F_X' = f_X$ \\
  media di $g(X)$ & $\sum_x g(x)\, p_X(x)$ & $\int g(x)\, f_X(x) \dd x$ \\
  funzione $Y = g(X)$ & $\sum_{g(x) = y} p_X(x)$ & $\sum_{g(x_i) = y} \frac{f_X(x_i)}{\abs{\dot g(x_i)}}$ \\
  varianza & $E[X^2] - \mu_X^2$ & $E[X^2] - \mu_X^2$ \\
  \bottomrule
\end{tabular}

Quasi ogni formula del discreto vale nel continuo con l'integrale al posto della somma e $f_X(x) \dd x$ al posto di
$p_X(x)$. L'unica davvero nuova è la densità di una funzione, con la divisione per $\abs{\dot g}$.
\end{scelta}

\scheda{valoremedioc}{
  \usoquando{Si vuole la media di $g(X)$ (un quadrato, una potenza, un costo): non serve la densità di $g(X)$.}
  \usocome{Si integra $g(x)$ pesata con la densità di $X$, sul supporto di $X$.}
  \usoesempio{$X \sim \mathcal{U}(0, 1)$ e $Y = X^2$: $E[Y] = \int_0^1 x^2 \dd x = \frac{1}{3}$. Passando dalla
    densità di $Y$, $\frac{1}{2\sqrt{y}}$ su $(0, 1)$, si trova lo stesso risultato con più fatica.}
}

\scheda{varianzac}{
  \usoquando{Si vuole la dispersione di una variabile continua. Per le variabili notevoli ci sono le formule pronte
    (\vv{uniformec}, \vv{esponenziale}, \vv{laplacevar}).}
  \usocome{Come nel discreto: $E[X^2]$ con il teorema fondamentale, poi si toglie $\mu_X^2$.}
  \usoesempio{$f_X(x) = 2x$ su $[0, 1]$: $E[X] = \frac{2}{3}$, $E[X^2] = \int_0^1 2x^3 \dd x = \frac{1}{2}$, quindi
    $\sigma_X^2 = \frac{1}{2} - \frac{4}{9} = \frac{1}{18}$.}
}

\scheda{quantizzazione}{
  \usoquando{Si quantizza un segnale con passo $\Delta$ e si vuole la potenza dell'errore (che ha media nulla, quindi
    potenza = varianza).}
  \usoattenzione{Vale se gli intervalli sono stretti rispetto alle variazioni della densità: allora l'errore è circa
    uniforme in ogni intervallo.}
  \usoesempio{Segnale in $[-1, 1]$ volt con 8 bit: $\Delta = \frac{2}{2^8} \simeq 0{,}0078$ V e $\sigma^2 =
    \frac{\Delta^2}{12} \simeq 5{,}09 \cdot 10^{-6}$ V$^2$; con 9 bit la potenza dell'errore è 4 volte più piccola.}
}

\scheda{dizionario}{
  \usoquando{Ogni volta che si conosce la formula nel discreto e serve quella nel continuo (o viceversa).}
  \usoattenzione{Per la funzione di una variabile continua non basta sommare le densità: va divisa ognuna per
    $\abs{\dot g}$ (\vv{nonmonotona}). Per il condizionamento si passa dalla CDF, perché $P(X = x) = 0$.}
  \usoesempio{Chebyshev vale anche nel continuo, ma spesso è largo: per $X \sim \mathcal{E}(1)$ ($\mu = \sigma = 1$),
    $P(\abs{X - 1} \geq 2) = P(X \geq 3) = e^{-3} \simeq 0{,}050$, contro il limite $\frac{1}{4}$.}
}
